% ID: FIS-FLU-PAS-001
% Titolo: Rapporto tra le superfici di un elevatore idraulico
% Disciplina: Fisica
% Area: Fluidostatica
% Argomento: Principio di Pascal
% Sottoargomento: Torchio idraulico
% Classe: Seconda liceo scientifico
% Difficolta: 3
% Tipo: problema_numerico
% Risultato: soluzione da aggiungere
% Tempo_stimato: 8 min
% Competenze: modellizzazione, proporzionalita, applicazione del principio di Pascal
% Prerequisiti: pressione, principio di Pascal, forza peso
% Tag: fisica, fluidostatica, principio_pascal, torchio_idraulico, pressione
% Fonte: verifica_fisica_fluidostatica_anonimizzata
% Autore: Riccardo Carli
% Licenza: CC-BY-SA-4.0

\begin{esercizio}
Un elevatore idraulico viene utilizzato per sollevare un'automobile all'interno
di un'officina. Se la massa della macchina e $3800\,\mathrm{kg}$ e la forza
che si esercita sulla superficie del pistone minore del torchio e
$140\,\mathrm{N}$, determina quante volte la superficie del secondo pistone e
piu grande di quella del primo.
\end{esercizio}

\begin{soluzione}
% Soluzione da aggiungere.
\end{soluzione}

